% \input{\pSections "ssec-korobeinikov"}

%%    %%    %%    %%    %%    %%    %%    %%    %%    %%
\subsection{Korobe\u{i}nikov}

\begin{frame}\frametitle{Problems in the Theory of Point Explosion in Gases}
\center
	\href{https://bookstore.ams.org/steklo-119}{
	\begin{overpic}[ scale = 0.2 ]
		{\pLocalGraphics cvr-korobinikov}
		%\put(50, 50)	{\colorbox{white}{$a+b$}}
	\end{overpic}}
	\\[10pt]
	\cite{korobeuinikov1976problems}
\end{frame}

%\begin{frame}\frametitle{Invariant Solutions}
%	%
%\small{
%\begin{equation}
%\begin{array}{rcl crc lcr clc rcl}
%		%
%	v &=& r 				&& \rho &=& e^{\beta t}R(r) 			&& p &=& e^{\beta t}P(r) \\[5pt]
%		%
%	v &=& r + e^{t}u(\lambda) 	&& \rho &=& e^{(\beta-2) t}R(\lambda) 	&& p &=& e^{\beta t}P(r)	&& \lambda &=&r e^{-t}\\[5pt]
%		%
%	v &=& \frac{1}{t}u(r) 	&& \rho &=& e^{2(\beta+1)}R(\lambda) 		&& p &=& e^{2^{t}}P(r) \\[5pt]
%		%
%	v &=& r + e^{t}u(\lambda) 	&& \rho &=& e^{(\beta-2) t}R(\lambda) 	&& p &=& e^{\beta t}P(r)	&& \lambda &=&r e^{-t}\\
%		%
%\end{array}
%\label{eq:invariant solutions}
%\end{equation}
%}
%	%
%\cite[p. 51, 1.134-137]{korobeuinikov1976problems}\\
%\cite{ovsjannikov1962group}
%\end{frame}

\begin{frame}\frametitle{Invariant Solutions: Cataloged}
	%
One-dimensional motions with \bl{spherical} or \bl{cylindrical} symmetry have \rd{only} the \bl{invariant solutions}:
\small{
\begin{equation}
\begin{array}{clclclcl}
		%
	&v = r 				&& \rho = e^{\beta t}R(r) 					&& p = e^{\beta t}P(r) \\[8pt]
		%
	&v = r + e^{t}u(\lambda) 	&& \rho = e^{(\beta-2) t}R(\lambda) 			&& p = e^{\beta t}P(\lambda)	&& \lambda = r e^{-t}\\[8pt]
		%
	&v = \frac{1}{t}u(r) 		&& \rho = e^{2(\beta+1)}R(\lambda) 			&& p = t^{2^{\beta}}P(r) \\[8pt]
		%
	&v =\frac{r}{t}u(\lambda) 	&& \rho = e^{(\beta+2\alpha-2) t}R(\lambda) 	&& p = r^{\beta}P(\lambda)	&& \lambda = tr^{-\alpha}\\
		%
\end{array}
\label{eq:invariant solutions}
\end{equation}
}
	%
\\[20pt]
\cite[p. 51, 1.134-137]{korobeuinikov1976problems} \\
\cite{ovsjannikov1962group}
\end{frame}
%
\begin{frame}\frametitle{Sedov Solutions}
Sedov's class of \bl{self-similar solutions} of 1D gas dynamics
		%
\begin{equation}
\begin{split}
	v &= \frac{r}{t}		V(\lambda), \\
	\rho &= \frac{a}{r^{k+3}t^{s}}	R(\lambda), \\
	p &= \frac{a}{r^{k+1}t^{s+2}}	P(\lambda), \\
	\lambda &= \frac{r}{b t^{d}}
\end{split}
\label{eq:sedov}
\end{equation}
		%
\ \\
Dimensional constants: $a$, $b$, \\
Numbers: $k$, $s$, $\delta$
\end{frame}
%
% https://tex.stackexchange.com/questions/144967/dividing-a-proof-into-several-frames
\begin{frame}\frametitle{Invariant Solutions of Gas Dynamics}
% Theorem 1, p. 52
	%
	\input{\pLocalTheorems "thm-invariant-solns-gasdynamics"}	%
\end{frame}
%
\begin{frame}\frametitle{Invariant Solutions of Gas Dynamics: Proof Outline}
% Theorem 1, p. 52
Let $t\to t + t_{0}$ where
	%
\begin{equation}
\begin{split}
	b &= \delta \tau \\
	V &= \delta \tilde{V} \\
	R &= \delta^{s} \tilde{R} \\
	P &= \delta^{s+2} \tilde{P}
\end{split}
\label{eq:proof-subs}
\end{equation}
\end{frame}
%
%
\begin{frame}\frametitle{Invariant Solutions of Gas Dynamics: Proof Outline}
% Theorem 1, p. 52
\begin{equation}
\begin{split}
	v &= \frac{r}{t}		\tilde{V}(\lambda), \\
	\rho &= \frac{a}{r^{k+3}t^{s}}	\tilde{R}(\lambda), \\
	p &= \frac{a}{r^{k+1}t^{s+2}}	\tilde{P}(\lambda), \\
	\lambda &= \frac{r}{r_{0}}e^{-t/\tau}
\end{split}
\label{eq:sedov-like}
\end{equation}
\ \\
Very much like \eqref{eq:sedov}
\end{frame}

%
\begin{frame}\frametitle{Invariant Solutions of Gas Dynamics: Comparison}
	%
\begin{table}[htp]
%\caption{default}
\begin{center}
\begin{tabular}{rll}
		%
	& Sedov \eqref{eq:sedov} & Proof \eqref{eq:sedov-like} \\\hline
	\\[-9pt]
		%
	$v$ 			& $\frac{r}{t} V (\lambda)$ 				&  $\frac{r}{t} \bl{\tilde{V}} (\lambda)$ \\[5pt]
	$\rho$ 		& $\frac{a}{r^{k+3}t^{s}} R (\lambda)$	& $\frac{a}{r^{k+3}t^{s}} \bl{\tilde{R}} (\lambda)$ \\[5pt]
	$p$ 			& $\frac{a}{r^{k+1}t^{s+2}} P(\lambda)$	& $\frac{a}{r^{k+1}t^{s+2}} \bl{\tilde{P}}(\lambda)$ \\[5pt]
	$\lambda$ 	& $\frac{r}{b t^{d}}$					& $\frac{r}{\bl{r_{0}}}\bl{e^{-t/\tau}}$
		%
\end{tabular}
\end{center}
\label{tab:proof-compare}
\end{table}%
	%
\end{frame}
%


\endinput  %  ==  ==  ==  ==  ==  ==  ==  ==  ==
